\section{Гармонические вложения: точные редукции и конечные семейства}

Во всём этом разделе $k$ --- поле характеристики нуль. Для согласования
с формулами репозитория используем моническую нормировку
\[
 C_n(t)=2T_n(t/2),\qquad C_0=2,\quad C_1=t,\quad
 C_{n+1}=tC_n-C_{n-1}.
\]
Она связана со стандартной нормировкой $T_n(\cos\theta)=\cos(n\theta)$
линейной заменой параметра и масштабированием координат. Поэтому вопрос
выпрямляемости одинаков для обеих нормировок. Обозначим
$\gamma_{a,b,c}=(C_a,C_b,C_c)$.

\paragraph{Что уже известно в литературе.}
В статье G.~Freudenburg и J.~Freudenburg,
\emph{Curves Defined by Chebyshev Polynomials},
\href{https://arxiv.org/abs/0902.3440}{arXiv:0902.3440},
предложение 4.1, установлено: если $a,b,c\geq2$, то $\gamma_{a,b,c}$
является замкнутым вложением тогда и только тогда, когда индексы попарно
взаимно просты. При наличии индекса $1$ вложение и его выпрямление
очевидны, даже без взаимной простоты остальных индексов.
Предложение 4.2 той же статьи даёт алгебраическую редукцию индексов;
следствия 4.1 и 4.2 устанавливают выпрямляемость всех вложений
$H(2,b,c)$ и $H(a,b,a+b)$ над $k$.

Работы P.-V.~Koseleff и D.~Pecker
\href{https://arxiv.org/abs/0909.3281}{\emph{Chebyshev diagrams for two-bridge knots}}
и \href{https://arxiv.org/abs/1203.4376}{\emph{Harmonic knots}}
дают топологическую классификацию семейств с первым индексом не больше
$4$. Эта классификация сама по себе не является классификацией
вложений по полиномиальным автоморфизмам над $\mathbb C$.

\paragraph{Явный полином, восстанавливающий параметр.}
Основные тождества имеют вид
\[
 C_mC_n=C_{m+n}+C_{|m-n|},\qquad C_m\circ C_n=C_{mn}
 \quad(m,n\geq1).
\]
Пусть $a$ нечётно и $a,b,c$ попарно взаимно просты. Существуют
положительные целые $r,s,q$ такие, что
\[
 |rb-sc|=1,\qquad rb+sc=qa.
\]
Действительно, после выбора положительного решения $rb-sc=1$ можно
заменить $(r,s)$ на $(r+cj,s+bj)$; второе условие сводится к разрешимой
конгруэнции с коэффициентом $2bc$, обратимым по модулю $a$.
Тогда
\[
 \boxed{\;C_r(C_b(t))C_s(C_c(t))-C_q(C_a(t))=t.\;}
\]
Это сертификат равенства $k[C_a,C_b,C_c]=k[t]$. Он удостоверяет
\emph{вложение}, но не утверждает, что восстанавливающий полином
$C_r(y)C_s(z)-C_q(x)$ является координатой кольца $k[x,y,z]$.

\paragraph{Элементарная редукция индексов.}
Если $c=ra+sb$, $r,s>0$, то
\[
 (x,y,z)\longmapsto(x,y,C_r(x)C_s(y)-z)
\]
--- полиномиальная инволюция, переводящая
$\gamma_{a,b,c}$ в $\gamma_{a,b,|ra-sb|}$.
В стандартной нормировке $T_n$ перед произведением нужен множитель $2$.
В HTML-тексте доказательства предложения 4.2 статьи Freudenburg--Freudenburg
этот множитель в записи инволюции пропущен; предшествующее тождество
напечатано с правильным множителем.

Сумма индексов при редукции строго убывает. После перестановок координат
процесс заканчивается либо индексом $1$, либо тройкой
\[
 2\leq A<B<C,\qquad C\notin\langle A,B\rangle,
 \qquad C\leq AB-A-B.
\]
Редукция обеспечивает конечный список кандидатов при фиксированной
исходной паре индексов. Она не утверждает, что различные конечные
тройки алгебраически неэквивалентны, и не даёт выпрямления всех таких троек.

\subsection{Алгебраические сдвиги третьего индекса}

Ниже приведён самостоятельный элементарный вывод из описания узлов
плоской чебышёвской кривой. Заявления о новизне этого следствия не делается.
Он полезен тем, что переводит определённые сравнения индексов в
\emph{конкретные полиномиальные сдвиги} третьей координаты.

\paragraph{Лемма о кольце узловой проекции.}
Пусть $a,b\geq2$, $\gcd(a,b)=1$ и
\[
 A=k[C_a(t),C_b(t)]\subset k[t],\qquad
 \delta=\frac{(a-1)(b-1)}2.
\]
Над алгебраическим замыканием плоская кривая
$(C_a(t),C_b(t))$ имеет ровно $\delta$ обыкновенных двойных точек
и не имеет других особенностей (предложение 3.1 статьи
Freudenburg--Freudenburg). Если $(s_\nu,t_\nu)$ --- две прообразные
точки каждого узла, то
\[
 \boxed{\quad A=\{f\in k[t]:f(s_\nu)=f(t_\nu)\text{ для всех }\nu\}.\quad}
\]
Равенства понимаются после расширения поля, если узлы не расщепляются над $k$.

\emph{Доказательство.}
Включение слева направо очевидно. Полиномы
$C_a^rC_b^s$, $r\geq0$, $0\leq s<a$, образуют базис $A$ над $k$:
порождаемость следует из соотношения $C_a(C_b)=C_b(C_a)$, монического
по $C_b$ степени $a$. Их ведущие степени $ra+sb$ попарно различны.
Удаление ведущих членов показывает, что коразмерность $A$ в $k[t]$
равна числу пробелов полугруппы $\langle a,b\rangle$, то есть $\delta$.
С другой стороны, $\delta$ равенств значений независимы: все $2\delta$
прообразных точек узлов различны, и на них применима обычная полиномиальная
интерполяция. Поэтому правая часть тоже имеет коразмерность $\delta$.
Это доказывает равенство над алгебраическим замыканием, а алгоритм
удаления ведущих членов даёт спуск к $k$.\hfill$\square$

\paragraph{Теорема о сдвиге.}
Пусть целые $c,c'\geq1$ удовлетворяют
\[
 c'\equiv\varepsilon_a c\pmod {2a},\qquad
 c'\equiv\varepsilon_b c\pmod {2b},\qquad
 \varepsilon_a,\varepsilon_b\in\{1,-1\}.
\]
Положим $\sigma=\varepsilon_a\varepsilon_b$. Тогда существует
$P\in\mathbb Q[x,y]$ такое, что
\[
 \boxed{\quad C_{c'}(t)-\sigma C_c(t)=P(C_a(t),C_b(t)).\quad}
\]
Следовательно,
\[
 S(x,y,z)=(x,y,\sigma z+P(x,y))
\]
--- полиномиальный автоморфизм над $\mathbb Q$, переводящий
$\gamma_{a,b,c}$ в $\gamma_{a,b,c'}$; его обратный имеет вид
$(x,y,\sigma(z-P(x,y)))$.

\emph{Доказательство.}
Прообразные точки каждого узла можно записать в виде
\[
 s=2\cos\pi\left(\frac{k}{a}+\frac{\ell}{b}\right),\qquad
 t=2\cos\pi\left(\frac{k}{a}-\frac{\ell}{b}\right).
\]
По формуле разности косинусов
\[
 C_c(s)-C_c(t)
 =-4\sin\frac{ck\pi}{a}\sin\frac{c\ell\pi}{b}.
\]
При замене $c$ на $c'$ правая часть умножается на $\sigma$.
По лемме $C_{c'}-\sigma C_c\in A$. Все исходные коэффициенты
рациональны, поэтому алгоритм членения по ведущей степени даёт
$P\in\mathbb Q[x,y]$.\hfill$\square$

В частности,
\[
 H(a,b,c+2ab)\simeq_{\mathrm{alg},\mathbb Q}H(a,b,c).
\]
Это равенство алгебраических классов, а не только равенство типов
вещественных узлов.

\paragraph{Алгоритм нахождения $P$.}
В монической нормировке достаточно повторять следующую операцию.
Для старшего члена $\lambda t^d$ вычислить единственное
$s\in\{0,\ldots,a-1\}$, удовлетворяющее $sb\equiv d\pmod a$,
и положить $r=(d-sb)/a$. Если $r\geq0$, вычесть
$\lambda C_a^rC_b^s$ и добавить $\lambda x^ry^s$ к $P$.
Иначе степень $d$ --- пробел полугруппы; её член переносится в остаток.
Для полиномов из теоремы остаток обязан быть нулевым.
Это даёт одновременно алгоритм, тождество-сертификат и его прямую проверку.

\paragraph{Число представителей при фиксированной паре.}
Среди классов $c\pmod {2ab}$ ровно $2\varphi(a)\varphi(b)$ удовлетворяют
$\gcd(c,ab)=1$. Четыре независимых выбора знаков в теореме действуют
на этом множестве свободно: совпадение двух знаков по модулю $2a$
или $2b$ означало бы соответственно $a\mid c$ или $b\mid c$.
Поэтому число нужных представителей не превосходит
\[
 \boxed{\quad \frac{\varphi(a)\varphi(b)}2.\quad}
\]
Это верхняя оценка числа алгебраических классов. Некоторые представители
могут быть эквивалентны более общими автоморфизмами.

\subsection{Два полных бесконечных семейства как следствие базовых моделей}

Для $(a,b)=(3,4)$ достаточно двух представителей. В таблице
третья колонка означает остатки индекса $c$, а не значения полиномов.
\[
\begin{array}{c|c|l}
 (a,b)&\text{представитель}&c\pmod{2ab}\\\hline
 (3,4)&1&1,7,17,23\\
 (3,4)&5&5,11,13,19
\end{array}
\]
Таким образом, подтверждённое выпрямление $H(3,4,5)$ над полем $K$
немедленно даёт выпрямление \emph{всех} $H(3,4,c)$,
$\gcd(c,12)=1$, над тем же $K$.

Для $(a,b)=(3,5)$ достаточно четырёх представителей:
\[
\begin{array}{c|c|l}
 (a,b)&\text{представитель}&c\pmod{2ab}\\\hline
 (3,5)&1&1,11,19,29\\
 (3,5)&2&2,8,22,28\\
 (3,5)&4&4,14,16,26\\
 (3,5)&7&7,13,17,23
\end{array}
\]
Первые две строки выпрямляются над $\mathbb Q$ по уже известным результатам;
третья после перестановки координат есть $H(3,4,5)$;
четвёртая есть $H(3,5,7)$.
Поэтому подтверждённые выпрямления двух базовых моделей над $K$
дают выпрямление \emph{всех} $H(3,5,c)$,
$\gcd(c,15)=1$, над тем же $K$.

Например, следующие тождества проверяются точным раскрытием:
\[
\begin{aligned}
 C_7+C_1&=C_3C_4,\\
 C_{11}+C_5&=C_3(C_4^2-2),\\
 C_{11}+C_1&=(C_3^2-2)C_5,\\
 C_{14}+C_4&=(C_3^3-3C_3)C_5,\\
 C_{13}+C_7&=C_3(C_5^2-2).
\end{aligned}
\]
Здесь $C_j$ справа в скобках обозначает значение $C_j(t)$,
а не композицию, кроме явно выписанных обычных операций.

\paragraph{Вычислительные сертификаты.}
Файл \texttt{periodic\_shears.py} строит $P$ точной рациональной
арифметикой и непосредственно проверяет подстановкой 191 совместимое
сравнение индексов для пар
$(2,3),(3,4),(3,5),(4,5),(3,7),(5,7)$ при
$1\leq c\leq20$, $c<c'\leq40$.
Кроме того, проверены девять отдельных примеров, включая
$(a,b,c,c')=(7,9,11,137)$.
Результаты и явные полиномы сохранены в
\texttt{periodic\_shears\_results.json}.
Эти вычисления проверяют реализацию теоремы, но доказательство
теоремы не основано на конечном переборе.

\paragraph{Граница результата.}
Для всех гармонических вложений получены критерий вложения,
явное восстановление параметра и конечная редукция при фиксированной
паре индексов. Из этого не следует выпрямляемость произвольного
$H(a,b,c)$. В частности, тип узла
$H(5,6,7)$ совпадает с типом $H(4,5,7)$ по теореме 4.4
Koseleff--Pecker, но это топологическая изотопия и не даёт
полиномиального автоморфизма между комплексными вложениями.
Не следует также переносить утверждения о $H(3,5,7)$ на
$(t^3-3t,t^5-5t^3+4t,t^7-7t^5+14t^3-7t)$:
средняя координата этой последней модели отличается от $C_5$.

\subsection{Следствие о минимальной степени поля определения}

Пусть $\gamma$ определена над $\mathbb Q$ и её вещественный образ
задаёт нетривиальный длинный узел. Если полиномиальное выпрямление
$\gamma$ вместе с обратным отображением определено над числовым полем
$L$, то $L$ не имеет вещественных вложений.
В самом деле, применение любого вложения $L\hookrightarrow\mathbb R$
ко всем коэффициентам сохранило бы тождества обратимости и
$\Phi(\gamma(t))=(t,0,0)$, дав невозможное вещественное выпрямление.
Полиномиальный автоморфизм $\mathbb R^3$ является собственным
гомеоморфизмом и продолжается на одноточечную компактификацию,
поэтому не переводит нетривиальный узел в прямую.

Следовательно, $L$ вполне мнимо и $[L:\mathbb Q]$ чётно.
Если выпрямление построено над $\mathbb Q(\sqrt{-3})$, степень $2$
минимальна среди возможных степеней числовых полей определения.
Это утверждение не означает минимальности степеней компонент
полиномиального автоморфизма и не выделяет единственное допустимое
квадратичное поле.

\subsection{Проверенное препятствие для естественного переноса $H(3,7,8)$}

Следующее тождество подсказывает очень короткую бирациональную редукцию:
\[
 (C_3^2-2)C_8-C_7^2+2=C_2.
\]
Поэтому отображение
\[
 \beta(x,y,z)=\bigl(x,y,(x^2-2)z-y^2+2\bigr)
\]
переводит $H(3,7,8)$ в заведомо выпрямляемую гармоническую модель
$H(3,7,2)$. Однако это \emph{не} доказывает выпрямляемости исходной
модели: $\beta^{-1}$ делит на $x^2-2$.

Более того, для этого конкретного $\beta$ условие графового центра,
используемое в принципе переноса репозитория с сохранением параметра,
выполнить невозможно. Действительно, пусть $r^2=2$. На слое $x=r$
центр есть парабола $Z=2-Y^2$. Если полиномиальный автоморфизм
$A_0=(x,H,W)$ переводит её в график $W=\varphi(H)$, то ограничение
$H$ на параболу обязано быть автоморфизмом аффинной прямой, то есть
\[
 H(r,Y,2-Y^2)=\lambda Y+\mu,\qquad\lambda\ne0.
\]
Если одновременно $H(C_3(t),C_7(t),C_2(t))=t$, то на трёх
различных корнях $C_3(t)=r$ должно выполняться
$t=\lambda C_7(t)+\mu$. Но точное деление даёт
\[
 C_7(t)\equiv-r t^2+(r^2-1)t+2r
 =-r t^2+t+2r\pmod{C_3(t)-r}.
\]
Сравнение коэффициентов квадратичного остатка требует
$\lambda r=0$, что невозможно. То же рассуждение работает, если
восстанавливаемый параметр есть $\alpha t+\eta$, $\alpha\ne0$.

Следовательно, именно эта бирациональная редукция не допускает
нужного графового переноса. Это препятствие к \emph{данному способу
выпрямления}, а не доказательство невыпрямляемости $H(3,7,8)$.

\paragraph{Общий фильтр знаменателей для кубической первой координаты.}
Тот же аргумент применим к любой модификации
\[
 \beta=(x,y,m(x)z+F(x,y))
\]
на исходной модели с первыми двумя координатами $(C_3,C_b)$.
Если существует исправляющий автоморфизм $(x,H,W)$, у которого
$H(\beta\gamma(t))=t$ и каждый приведённый центр становится графиком
над $H$, то для каждого корня $r$ полинома $m$ необходимо
\[
 t=\lambda C_b(t)+\mu\quad\text{в }k[t]/(C_3(t)-r),\qquad\lambda\ne0.
\]
Этот вывод остаётся верен и для кратных корней $C_3(t)-r$:
он получается подстановкой в полиномиальное тождество по модулю
$C_3(t)-r$, а не только проверкой геометрических точек.

Положим $V_n(x)=U_n(x/2)$, $V_{-1}=0$, $V_0=1$,
$V_{n+1}=xV_n-V_{n-1}$. По индукции из чебышёвской рекурсии
получаем точные остатки
\[
\begin{aligned}
 C_{3n+1}(t)&\equiv -V_{n-1}(x)t^2+V_n(x)t+2V_{n-1}(x),\\
 C_{3n+2}(t)&\equiv V_n(x)t^2-V_{n-1}(x)t-2V_n(x)
 \pmod{C_3(t)-x}.
\end{aligned}
\]
Следовательно, приведённый многочлен $m_{\mathrm{red}}$ обязан делить
$V_{n-1}$ при $b=3n+1$ и $V_n$ при $b=3n+2$.
Соседние $V_n,V_{n-1}$ взаимно просты, так что линейный коэффициент
в остатке на разрешённых корнях не обнуляется.
Например, при $b=5$ и $b=7$ допустимы только корни $m$ в $x=0$;
при $b=8$ --- только $x=\pm1$; при $b=11$ --- только
$x=0,\pm\sqrt2$.
Это необходимое условие для указанного механизма переноса;
достаточность и условия на кратности корней здесь не установлены.
